% Options for packages loaded elsewhere
\PassOptionsToPackage{unicode}{hyperref}
\PassOptionsToPackage{hyphens}{url}
\PassOptionsToPackage{dvipsnames,svgnames,x11names}{xcolor}
\documentclass[
  10pt,
]{article}
\usepackage{xcolor}
\usepackage[margin=2.4cm]{geometry}
\usepackage{amsmath,amssymb}
\setcounter{secnumdepth}{-\maxdimen} % remove section numbering
\usepackage{iftex}
\ifPDFTeX
  \usepackage[T1]{fontenc}
  \usepackage[utf8]{inputenc}
  \usepackage{textcomp} % provide euro and other symbols
\else % if luatex or xetex
  \usepackage{unicode-math} % this also loads fontspec
  \defaultfontfeatures{Scale=MatchLowercase}
  \defaultfontfeatures[\rmfamily]{Ligatures=TeX,Scale=1}
\fi
\usepackage{lmodern}
\ifPDFTeX\else
  % xetex/luatex font selection
  \setmainfont[]{Libertinus Serif}
  \setmonofont[Scale=0.85]{Latin Modern Mono}
  \setmathfont[]{Latin Modern Math}
\fi
% Use upquote if available, for straight quotes in verbatim environments
\IfFileExists{upquote.sty}{\usepackage{upquote}}{}
\IfFileExists{microtype.sty}{% use microtype if available
  \usepackage[]{microtype}
  \UseMicrotypeSet[protrusion]{basicmath} % disable protrusion for tt fonts
}{}
\makeatletter
\@ifundefined{KOMAClassName}{% if non-KOMA class
  \IfFileExists{parskip.sty}{%
    \usepackage{parskip}
  }{% else
    \setlength{\parindent}{0pt}
    \setlength{\parskip}{6pt plus 2pt minus 1pt}}
}{% if KOMA class
  \KOMAoptions{parskip=half}}
\makeatother
\usepackage{longtable,booktabs,array}
\newcounter{none} % for unnumbered tables
\usepackage{calc} % for calculating minipage widths
% Correct order of tables after \paragraph or \subparagraph
\usepackage{etoolbox}
\makeatletter
\patchcmd\longtable{\par}{\if@noskipsec\mbox{}\fi\par}{}{}
\makeatother
% Allow footnotes in longtable head/foot
\IfFileExists{footnotehyper.sty}{\usepackage{footnotehyper}}{\usepackage{footnote}}
\makesavenoteenv{longtable}
\setlength{\emergencystretch}{3em} % prevent overfull lines
\providecommand{\tightlist}{%
  \setlength{\itemsep}{0pt}\setlength{\parskip}{0pt}}

\providecommand{\E}{\mathbb{E}}
\providecommand{\N}{\mathbb{N}}
\providecommand{\Z}{\mathbb{Z}}
\providecommand{\R}{\mathbb{R}}
\providecommand{\eps}{\varepsilon}
\providecommand{\ceil}[1]{\left\lceil #1\right\rceil}
\providecommand{\floor}[1]{\left\lfloor #1\right\rfloor}
\AtBeginDocument{\let\setminus\smallsetminus}
\usepackage[most]{tcolorbox}
\newtcolorbox{warning}{colback=red!5,colframe=red!65!black,boxrule=0.8pt,arc=2pt,left=6pt,right=6pt,top=5pt,bottom=5pt}
\usepackage{etoolbox}
\AtBeginEnvironment{longtable}{\small}
\setlength{\emergencystretch}{3em}
\usepackage{fancyhdr}
\pagestyle{fancy}
\fancyhf{}
\fancyhead[L]{\small\bfseries UNREFEREED GPT-6 ASTRA OUTPUT}
\fancyhead[R]{\small\thepage}
\renewcommand{\headrulewidth}{0.4pt}
\usepackage{bookmark}
\IfFileExists{xurl.sty}{\usepackage{xurl}}{} % add URL line breaks if available
\urlstyle{same}
\hypersetup{
  pdftitle={Finding k-edge-outerplanar graph embeddings},
  pdfauthor={Writeup: gpt-6-astra (single model pass)},
  colorlinks=true,
  linkcolor={blue!50!black},
  filecolor={Maroon},
  citecolor={Blue},
  urlcolor={blue!50!black},
  pdfcreator={LaTeX via pandoc}}

\title{Finding k-edge-outerplanar graph embeddings}
\usepackage{etoolbox}
\makeatletter
\providecommand{\subtitle}[1]{% add subtitle to \maketitle
  \apptocmd{\@title}{\par {\large #1 \par}}{}{}
}
\makeatother
\subtitle{Unrefereed candidate proof by GPT-6 Astra}
\author{Writeup: \texttt{gpt-6-astra} (single model pass)}
\date{Catalog id
\texttt{finding\_k\_edge\_outerplanar\_graph\_embeddings} --- generated
2026-09-16}

\begin{document}
\maketitle

\begin{warning}

\textbf{UNREFEREED MODEL OUTPUT.} This document was selected solely
because GPT-6 Astra labelled its own result
\texttt{would\_publish:\ true} and returned \texttt{proved} or
\texttt{disproved}. It has not passed the adversarial LLM referee used
for the earlier Sol campaign, has not been checked by a human
mathematician, and has not been checked for novelty. Treat every
mathematical and bibliographic claim below as unverified.

\end{warning}

\section{Summary and provenance}\label{summary-and-provenance}

{\def\LTcaptype{none} % do not increment counter
\begin{longtable}[]{@{}
  >{\raggedright\arraybackslash}p{(\linewidth - 2\tabcolsep) * \real{0.5000}}
  >{\raggedright\arraybackslash}p{(\linewidth - 2\tabcolsep) * \real{0.5000}}@{}}
\toprule\noalign{}
\endhead
\bottomrule\noalign{}
\endlastfoot
Catalog id &
\texttt{finding\_k\_edge\_outerplanar\_graph\_embeddings} \\
Catalog entry &
\href{https://graph-theory-ai.github.io/graph-conjectures/op/finding_k_edge_outerplanar_graph_embeddings/}{Finding
k-edge-outerplanar graph embeddings} \\
Source corpus & opg \\
Campaign leg & OpenProblemGarden first pass (\texttt{attacks\_opg}) \\
Source paper / entry &
http://www.openproblemgarden.org/op/finding\_k\_edge\_outerplanar\_graph\_embeddings \\
Model verdict & \textbf{proved} (confidence: high) \\
Model's one-line claim & A polynomial-time SPQR dynamic program, using a
two-bin scheduling recurrence at parallel nodes, computes a minimum
edge-outerplanar embedding. \\
Model & \texttt{gpt-6-astra}, reasoning effort \texttt{max},
\texttt{mode=pro}, flex service tier \\
Original artifact &
\texttt{attacks\_opg/finding\_k\_edge\_outerplanar\_graph\_embeddings/output.md} \\
Independent review & \textbf{None} \\
Model's caveats & Uses the standard SPQR decomposition theorem; novelty
and current literature status have not been independently checked. \\
\end{longtable}
}

\section{Problem statement}\label{problem-statement}

Conjecture It has been shown that a \(k\) -outerplanar embedding for
which \(k\) is minimal can be found in polynomial time. Does a similar
result hold for \(k\) -edge-outerplanar graphs?

\subsection{Catalog context}\label{catalog-context}

A \(k\) -outerplanar graph {[}Baker{]} with \(k > 0\) is a planar graph
having an embedding with at most \(k\) layers of vertices such that
after removing iteratively the vertices (and their adjacent edges) lying
on the outer face \(k\) times, we obtain the empty graph. A \(k\)
-edge-outerplanar graph {[}Bentz{]} is defined to be a planar graph
having an embedding with at most \(k\) layers of edges such that after
removing iteratively the edges lying on the outer face \(k\) times, we
obtain a graph with no edge. All \(k\) -edge-outerplanar graphs are
\(k\) -outerplanar graphs. Given a planar graph, Bienstock and Monma
have shown that a \(k\) -outerplanar embedding for which \(k\) is
minimal can be found in polynomial time. Does a similar result hold for
\(k\) -edge-outerplanar graphs?

\section{Astra writeup}\label{astra-writeup}

\emph{The text below is the model output verbatim, apart from moving its
machine-readable verdict block into the summary above and shifting
Markdown heading levels for this document. Mechanical TeX defects and
equation tags were normalized where needed for compilation.}

\subsection{A polynomial algorithm for minimum
edge-outerplanarity}\label{a-polynomial-algorithm-for-minimum-edge-outerplanarity}

I give an algorithmic proof. The main point is that the dual distance
across a pertinent SPQR component is independent of its embedding. The
remaining ordering problem at a parallel node has a polynomial dynamic
program.

The proof also handles cutvertices. They require some care: one cannot
simply take the maximum of the optimum values of the blocks, because an
attachment vertex need not lie on the chosen outer face of its block.

\subsubsection{1. Statement}\label{statement}

\textbf{Theorem.} Given a finite planar graph \(G\), one can compute, in
polynomial time, a plane embedding of \(G\) minimizing the number of
rounds of simultaneous deletion of all edges incident with the outer
face.

The algorithm returns a combinatorial embedding and its outer face. An
intentionally coarse bound for the implementation below is \(O(N^8)\)
operations on \(O(\log N)\)-bit integers, where \(N=|V(G)|+|E(G)|\).

For an edgeless graph, take the optimum to be \(0\); if the parameter is
required to be positive, report \(1\) instead.

The only standard structural input is the SPQR decomposition theorem for
biconnected planar graphs, including its description of all planar
embeddings. No result about optimizing outerplanarity is used.

\begin{center}\rule{0.5\linewidth}{0.5pt}\end{center}

\subsubsection{2. Edge peeling is an edge-distance objective in the
dual}\label{edge-peeling-is-an-edge-distance-objective-in-the-dual}

Let \(G\) be a plane graph, let \(D=G^*\) be its geometric dual, and let
\(r\) be the dual vertex corresponding to the outer face. Loops and
parallel edges are retained in the dual.

Write \[
d(f)=\operatorname{dist}_D(r,f).
\]

For a primal edge \(e\), let \(f_e,g_e\) be its two incident faces,
possibly equal.

\textbf{Lemma 1.} The round in which \(e\) is deleted is \[
\operatorname{layer}(e)
=
1+\min\{d(f_e),d(g_e)\}.
\qquad\text{(1)}
\]

\textbf{Proof.} After round \(t\), the original faces absorbed into the
outer face are exactly those at dual distance at most \(t\) from \(r\).
This follows inductively: deleting the current outer-face edges absorbs
precisely the neighboring original faces.

An edge is therefore exposed at the beginning of round \(t\) precisely
when one of its incident original faces has distance at most \(t-1\).
Its first such round is (1). This also covers a bridge, whose two
incident faces are equal. \(\square\)

Thus the edge-outerplanarity of this particular embedding is \[
K(G,r)
=
\max_{e\in E(G)}
\left(1+\min\{d(f_e),d(g_e)\}\right).
\qquad\text{(2)}
\]

In particular, the objective is \textbf{not quite the eccentricity of
\(r\) in the dual}. For example, in a tetrahedral embedding of \(K_4\),
the outer-face dual vertex has eccentricity \(1\), but edge peeling
takes \(2\) rounds.

We will also use the following immediate consequence.

\textbf{Deletion monotonicity.} Deleting edges from a plane graph, while
keeping the inherited outer face, cannot increase the deletion round of
any remaining edge.

Indeed, deletion merges faces, so dual distances cannot increase.

\begin{center}\rule{0.5\linewidth}{0.5pt}\end{center}

\subsubsection{3. A weighted problem needed for
cutvertices}\label{a-weighted-problem-needed-for-cutvertices}

For a plane graph \(B\), define \[
a_B(v)=\min\{d(f):f\text{ is incident with }v\}.
\qquad\text{(3)}
\] This is the distance from the outer face to a nearest face incident
with \(v\).

Give the vertices nonnegative integer weights \(w(v)\), and define \[
\Psi(B,w)
=
\max\left\{
K(B,r),\
\max_{v\in V(B)}\bigl(w(v)+a_B(v)\bigr)
\right\}.
\qquad\text{(4)}
\]

We first solve the following problem for a biconnected planar graph
\(B\):

\begin{itemize}
\tightlist
\item
  minimize \(\Psi(B,w)\) over its embeddings and outer faces;
\item
  optionally require a specified vertex \(c\) to be incident with the
  outer face.
\end{itemize}

The weights will eventually be optimum values of graphs attached at
cutvertices.

A zero weight imposes no additional restriction: for every nonisolated
vertex \(v\), an edge bounding a nearest incident face has layer
\(a_B(v)+1\), so \(a_B(v)<K(B,r)\).

\begin{center}\rule{0.5\linewidth}{0.5pt}\end{center}

\subsubsection{4. Rooted components and their
profiles}\label{rooted-components-and-their-profiles}

Fix an edge of a biconnected graph and root its SPQR decomposition at
that edge.

It is convenient to use the usual edge-substitution formulation. A
pertinent component is a two-terminal primal graph \(H\), with distinct
poles \(s,t\), together with an additional reference edge \(\rho=st\).
The augmented graph \[
\widehat H=H+\rho
\] is biconnected and planar.

For an embedding of \(\widehat H\):

\begin{itemize}
\tightlist
\item
  let \(A,B\) be the faces on the two sides of \(\rho\);
\item
  take the dual of \(\widehat H\);
\item
  delete the dual edge \(\rho^*\).
\end{itemize}

The resulting connected two-terminal dual graph is denoted \(D_H\), with
terminals \(A,B\). It has one edge for each real edge of \(H\).

\paragraph{4.1 The terminal distance is
embedding-invariant}\label{the-terminal-distance-is-embedding-invariant}

Set \[
\ell(H)=\operatorname{dist}_{D_H}(A,B).
\]

\textbf{Lemma 2.} The integer \(\ell(H)\) is independent of the
embedding choices in the rooted SPQR subtree. Moreover, \[
1\le \ell(H)\le |E(H)|.
\qquad\text{(5)}
\]

\textbf{Proof.} Use induction through the SPQR decomposition.

A real-edge leaf gives a single dual edge, so its value is \(1\).

At an internal node, replacing a primal skeleton edge by a child
component replaces the corresponding dual edge by the child's
two-terminal dual graph. Consequently, for distances between skeleton
faces, that child may be replaced by an edge of length \(\ell(H_i)\).

There are three skeleton types.

\begin{itemize}
\tightlist
\item
  For an \(S\)-node, the primal skeleton is a cycle. After removing the
  reference dual edge, its dual is a parallel composition. Hence \[
  \ell(H)=\min_i\ell(H_i).
  \]
\item
  For a \(P\)-node, the primal skeleton is a dipole. Its dual, after
  removing the reference dual edge, is a path through the children.
  Hence \[
  \ell(H)=\sum_i\ell(H_i),
  \] independently of their order.
\item
  For an \(R\)-node, the skeleton has a unique spherical embedding up to
  reflection. Its terminal distance is the shortest-path distance in the
  fixed dual skeleton with child-edge lengths \(\ell(H_i)\). Reflection
  does not change that distance.
\end{itemize}

This proves invariance. The bounds follow because \(D_H\) is connected
and has \(|E(H)|\) unit-length edges. \(\square\)

This invariance is what permits a polynomial number of states.

\paragraph{4.2 Profile definition}\label{profile-definition}

Weights on the poles \(s,t\) are temporarily ignored. All other vertex
weights are included.

For an integer \(d\) with \[
0\le d\le \ell(H),
\] give the two dual terminals external distance labels \(0,d\). In a
particular embedding, the resulting distance label of a dual vertex
\(z\) is \[
\delta(z)
=
\min\left\{
\operatorname{dist}_{D_H}(A,z),\
d+\operatorname{dist}_{D_H}(B,z)
\right\}.
\qquad\text{(6)}
\]

For that embedding, take the maximum of:

\begin{enumerate}
\def\labelenumi{\arabic{enumi}.}
\tightlist
\item
  for every real dual edge \(xy\), \[
  1+\min\{\delta(x),\delta(y)\};
  \]
\item
  for every internal primal vertex \(v\ne s,t\), \[
  w(v)+\min_{f\ni v}\delta(f).
  \]
\end{enumerate}

Let \(F_H(d)\) be the minimum of this maximum over the embeddings of the
pertinent component, including reflection.

Reflection exchanges the two boundary faces while preserving the primal
pole labels. Therefore, if the actual terminal labels are \(a,b\), with
\[
|a-b|\le \ell(H),
\] the optimum contribution of the component is \[
\min\{a,b\}+F_H(|a-b|).
\qquad\text{(7)}
\]

Each profile has at most \(|E(H)|+1\) entries.

For a real-edge leaf, \[
\ell(H)=1,\qquad F_H(0)=F_H(1)=1.
\qquad\text{(8)}
\]

\begin{center}\rule{0.5\linewidth}{0.5pt}\end{center}

\subsubsection{5. Composition over a fixed
skeleton}\label{composition-over-a-fixed-skeleton}

Consider a skeleton \(\Sigma\), with its reference edge, whose other
edges correspond to children \(H_i\). Write \[
\ell_i=\ell(H_i),\qquad F_i=F_{H_i}.
\]

Let \(W\) be the skeleton dual with the reference dual edge removed,
giving its child edge \(i\) length \(\ell_i\). Its reference-face
terminals are \(A,B\).

For prescribed terminal labels \(a,b\), compute \[
\delta(z)=
\min\left\{
a+\operatorname{dist}_W(A,z),\
b+\operatorname{dist}_W(B,z)
\right\}.
\qquad\text{(9)}
\]

If child edge \(i\) has dual endpoints \(x_i,y_i\), its optimum
contribution is \[
\min\{\delta(x_i),\delta(y_i)\}
+
F_i\bigl(|\delta(x_i)-\delta(y_i)|\bigr).
\qquad\text{(10)}
\] The profile argument is valid because \[
|\delta(x_i)-\delta(y_i)|\le \ell_i.
\]

\paragraph{Why child optimization is
independent}\label{why-child-optimization-is-independent}

Distances between skeleton faces depend only on the invariant lengths
\(\ell_i\), not on the child embeddings.

For an internal face of a child, every path from outside the child
enters through one of its two boundary faces. Thus its distance is
obtained from those two boundary labels exactly as in (6). Each child
can consequently be optimized independently for its prescribed labels.

\paragraph{Vertex weights localize to the
skeleton}\label{vertex-weights-localize-to-the-skeleton}

Suppose \(v\) is a skeleton vertex other than the parent poles. A new
face inside a child incident with \(v\) has distance at least the
minimum of the child's two boundary labels. Both boundary faces are
themselves incident with \(v\).

It follows that, after all substitutions, \[
\min_{f\ni v}\delta(f)
=
\min_{\substack{f\text{ a skeleton face}\\f\ni v}}\delta(f).
\qquad\text{(11)}
\]

Every internal primal vertex is either:

\begin{itemize}
\tightlist
\item
  internal to exactly one child; or
\item
  a skeleton vertex other than the parent poles.
\end{itemize}

Thus (10), together with the skeleton-vertex terms from (11), accounts
for every edge and every relevant vertex weight.

Define \[
C_\Sigma(a,b)
=
\max\left\{
\begin{array}{l}
\displaystyle
\max_i\left[
\min\{\delta(x_i),\delta(y_i)\}
+
F_i\bigl(|\delta(x_i)-\delta(y_i)|\bigr)
\right],
\\[2ex]
\displaystyle
\max_{v\in V(\Sigma)\setminus\{s,t\}}
\left[
w(v)+\min_{f\ni v}\delta(f)
\right]
\end{array}
\right\}.
\qquad\text{(12)}
\]

For a fixed skeleton embedding, this is the exact optimum over its child
embeddings.

\paragraph{\texorpdfstring{5.1 \(R\)-nodes}{5.1 R-nodes}}\label{r-nodes}

Compute \(\ell(H)\) as the weighted terminal distance in \(W\). Both
skeleton reflections must be considered: \[
F_H(d)=
\min\{C_\Sigma(0,d),\,C_\Sigma(d,0)\},
\qquad 0\le d\le \ell(H).
\qquad\text{(13)}
\]

The minimum in (13) is important: the two faces adjacent to a reference
edge need not have equivalent distance behavior.

\paragraph{\texorpdfstring{5.2 \(S\)-nodes}{5.2 S-nodes}}\label{s-nodes}

The two skeleton faces are common terminals of all children. Every
internal skeleton vertex is incident with both faces. Therefore, \[
\ell(H)=\min_i\ell_i
\] and \[
F_H(d)
=
\max\left\{
\max_i F_i(d),\
\max_{v\in V(\Sigma)\setminus\{s,t\}}w(v)
\right\}.
\qquad\text{(14)}
\]

It remains to treat \(P\)-nodes, where the children can be arbitrarily
ordered.

\begin{center}\rule{0.5\linewidth}{0.5pt}\end{center}

\subsubsection{\texorpdfstring{6. The polynomial recurrence at a
\(P\)-node}{6. The polynomial recurrence at a P-node}}\label{the-polynomial-recurrence-at-a-p-node}

At a \(P\)-node, the dual network is a series composition of the
children in an arbitrary order. There are no internal skeleton vertices,
so there are no additional skeleton-vertex weights.

Put \[
\Lambda=\sum_{i=1}^{r}\ell_i=\ell(H),
\qquad
p_i=F_i(\ell_i).
\qquad\text{(15)}
\]

The quantity \(p_i\) is the optimum one-sided cost of child \(i\).
Indeed, with terminal labels \(0,\ell_i\), the first terminal alone
determines all distances, since \[
\operatorname{dist}(A,z)
\le
\ell_i+\operatorname{dist}(B,z).
\]

Fix a profile argument \(d\in\{0,\ldots,\Lambda\}\). Along a weighted
series path, a skeleton vertex at distance \(x\) from the left terminal
has label \[
\delta(x)=\min\{x,\ d+\Lambda-x\}.
\qquad\text{(16)}
\]

The two expressions meet at \[
h=\frac{\Lambda+d}{2}.
\]

Choose a child \(j\) whose interval contains \(h\). If \(h\) is at a
junction, choose either adjacent child. Call \(j\) the central child.

Let:

\begin{itemize}
\tightlist
\item
  \(X\) be the total length of the children before \(j\);
\item
  \(Y\) be the total length of the children after \(j\).
\end{itemize}

Then \[
X+Y+\ell_j=\Lambda.
\] The labels at the two ends of the central child are \[
\alpha=X,\qquad \beta=d+Y.
\]

The condition that the meeting point lies in the central child is
exactly \[
|\alpha-\beta|\le \ell_j.
\qquad\text{(17)}
\]

Every child strictly before \(j\) is reached optimally from the left;
every child strictly after \(j\) is reached optimally from the right.
Their costs are consequently one-sided costs.

Thus, except for the central child, the problem is an assignment of
children to two arms, followed by ordering within each arm.

\paragraph{6.1 Optimal ordering within an
arm}\label{optimal-ordering-within-an-arm}

For child \(i\), define \[
q_i=p_i-\ell_i.
\qquad\text{(18)}
\]

On either arm, listed from its source toward the central child, there is
an optimum order of nonincreasing \(q_i\).

To verify this, suppose \(q_i\ge q_j\), and compare consecutive orders
\(ij\) and \(ji\), starting at accumulated label \(A\). The two costs in
order \(ij\) are \[
A+\ell_i+q_i,\qquad
A+\ell_i+\ell_j+q_j.
\] Both are at most \[
A+\ell_i+\ell_j+q_i,
\] which is one of the costs in order \(ji\). Swapping an inverted pair
therefore cannot worsen the maximum. Later costs are unchanged.

Reordering within an arm does not change its total length, so it also
leaves the central-child condition and contribution unchanged.

\paragraph{6.2 Assigning children to the two
arms}\label{assigning-children-to-the-two-arms}

Fix the central child \(j\), and sort the remaining children by
nonincreasing \(q_i\).

After processing some initial segment of this list, let its total length
be \(S\). Define \[
T(x)
\] to be the minimum possible maximum cost so far when the left arm has
total length \(x\), and the right arm has total length \(S-x\).

Initially, \[
T(0)=0,
\] with all other states infeasible.

For the next child \(i\), there are two transitions.

\begin{itemize}
\tightlist
\item
  Put it on the left: \[
  T_{\rm new}(x+\ell_i)
  \ \leftarrow\
  \min\left\{
  T_{\rm new}(x+\ell_i),\
  \max\{T(x),\,x+p_i\}
  \right\}.
  \qquad\text{(19)}
  \]
\item
  Put it on the right: \[
  T_{\rm new}(x)
  \ \leftarrow\
  \min\left\{
  T_{\rm new}(x),\
  \max\{T(x),\,d+(S-x)+p_i\}
  \right\}.
  \qquad\text{(20)}
  \]
\end{itemize}

Both arm orders are subsequences of the common sorted list. The actual
left-to-right series order is the left list, followed by \(j\), followed
by the reverse of the right list. Hence every assignment produced by
this dynamic program is realizable.

After all noncentral children have been processed, a state \(X\) has \[
Y=\Lambda-\ell_j-X.
\]

If \[
|X-(d+Y)|\le \ell_j,
\] its full cost is \[
\max\left\{
T(X),\
\min\{X,d+Y\}
+
F_j\bigl(|X-(d+Y)|\bigr)
\right\}.
\qquad\text{(21)}
\]

Take the minimum of (21) over all feasible \(X\), and then over all
choices of \(j\). The result is \(F_H(d)\).

\paragraph{6.3 Completeness of the
recurrence}\label{completeness-of-the-recurrence}

Every ordered series composition has a child containing the meeting
point in (16). All other children are one-sided. Sorting each arm as
above cannot worsen its contribution, and does not change the central
contribution.

Conversely, every state satisfying (17) gives a valid ordering with
precisely the terminal labels used in (21). Independently realizing the
appropriate child profiles realizes the stated cost.

Thus the recurrence is exact, including:

\begin{itemize}
\tightlist
\item
  a meeting point at a junction;
\item
  a meeting point at the right endpoint, when \(d=\Lambda\);
\item
  a composition with only one child.
\end{itemize}

\paragraph{6.4 Complexity}\label{complexity}

For fixed \(j,d\), the arm-assignment dynamic program has \(O(\Lambda)\)
states and \(r-1\) stages. Computing the whole profile takes \[
O(r^2\Lambda^2)
\qquad\text{(22)}
\] operations, apart from lower-order sorting work.

The lengths here are not arbitrary binary-encoded numbers:
\(\Lambda\le |E(H)|\). This is why the assignment dynamic program is
polynomial in the graph size.

\begin{center}\rule{0.5\linewidth}{0.5pt}\end{center}

\subsubsection{7. Solving the weighted biconnected
problem}\label{solving-the-weighted-biconnected-problem}

Let \(B\) be a biconnected planar graph with vertex weights \(w\). Fix a
real edge \(e=st\), and require \(e\) to be incident with the outer
face.

Root the SPQR decomposition at \(e\), and let \[
H=B-e,
\] using \(e\) as the reference edge.

In the full dual, the reference dual edge joins the outer-face terminal
to the other terminal by an edge of length \(1\). Thus the terminal
distance labels are \(0,1\).

The reference primal edge itself has layer \(1\). Its endpoints are
incident with the outer face, so their weighted contributions are
\(w(s),w(t)\). Consequently, the optimum with \(e\) exposed is \[
\operatorname{Opt}(B,w;e)
=
\max\{1,\ F_H(1),\ w(s),\ w(t)\}.
\qquad\text{(23)}
\]

All other edges and vertex weights are included in \(F_H(1)\).

It follows that \[
\operatorname{Opt}(B,w)
=
\min_{e\in E(B)}\operatorname{Opt}(B,w;e).
\qquad\text{(24)}
\]

If a specified vertex \(c\) must be incident with the outer face, use \[
\operatorname{Opt}(B,w;c)
=
\min_{\substack{e\in E(B)\\e\ni c}}
\operatorname{Opt}(B,w;e).
\qquad\text{(25)}
\]

Every outer-face incidence at \(c\) includes an incident edge, so (25)
loses no embedding.

Equations (8), (13), (14), and (19)--(21) provide a bottom-up polynomial
computation of all required profiles. The standard completeness property
of the SPQR decomposition ensures that all embeddings are covered.

\begin{center}\rule{0.5\linewidth}{0.5pt}\end{center}

\subsubsection{8. Cutvertices}\label{cutvertices}

The following gluing lemma supplies the reduction to weighted blocks.

\paragraph{8.1 Gluing lemma}\label{gluing-lemma}

Suppose a connected graph consists of a block \(B\), together with
connected branches \(H_{v,i}\) attached at vertices \(v\in V(B)\). Each
branch meets \(B\) only at \(v\); different branches are otherwise
disjoint. As in a rooted block-cut tree, each nontrivial branch minus
its attachment vertex is connected.

Let \[
h_{v,i}
\] be the minimum edge-outerplanarity of \(H_{v,i}\) among embeddings in
which \(v\) is incident with its outer face. Set \[
w(v)=\max_i h_{v,i},
\qquad\text{(26)}
\] with \(w(v)=0\) when there is no branch.

\textbf{Lemma 3.} Among embeddings of the whole graph whose outer face
is incident with an edge of \(B\), the optimum equals \[
\min_{\text{embeddings and outer faces of }B}\Psi(B,w).
\qquad\text{(27)}
\]

The analogous assertion holds with a required outer-face incidence in
\(B\).

\textbf{Proof: upper bound.} Choose an embedding of \(B\). For each
\(v\), choose an incident face \(f_v\) satisfying \[
d(f_v)=a_B(v).
\]

Embed every \(H_{v,i}\) in a small region of \(f_v\) adjoining \(v\),
using an optimum embedding with \(v\) on its outer face. These regions
can be chosen pairwise disjoint, except for their prescribed attachment
vertices.

In the dual, each branch dual is joined to the dual of \(B\) at the
single vertex \(f_v\). Distances between faces of \(B\) do not change.
All edge layers in the branch increase by exactly \(a_B(v)\).

The resulting value is therefore \[
\max\left\{
K(B,r),\
\max_{v,i}\bigl(a_B(v)+h_{v,i}\bigr)
\right\}
=
\Psi(B,w).
\]

\textbf{Proof: lower bound.} Consider any embedding of the whole graph
whose outer face is incident with an edge of \(B\). Restrict it to
\(B\), obtaining a fixed embedding of \(B\).

For one branch \(H_{v,i}\), delete all other branches. By deletion
monotonicity, this cannot increase any remaining layer.

The graph \(H_{v,i}-v\) lies in a single face \(f\) of \(B\), incident
with \(v\). Because an edge of \(B\) is on the outer face, the branch
does not enclose \(B\). In its induced standalone embedding, \(v\) is
therefore incident with the outer face.

The dual of \(B\cup H_{v,i}\) is obtained by identifying the outer-face
vertex of \(H_{v,i}^*\) with the vertex \(f\) of \(B^*\). Hence the
largest layer among the branch edges is at least \[
d_B(f)+h_{v,i}
\ge
a_B(v)+h_{v,i}.
\]

Also, restricting to \(B\) shows that the whole graph has value at least
\(K(B,r)\). Taking the maximum over branches proves the lower bound.
\(\square\)

Notice that this argument does not assume that an arbitrary embedding
has nonnested branches. Nesting is handled in the lower bound by
deleting all other branches. The upper bound constructs a nonnested
embedding attaining the weighted objective.

\paragraph{8.2 Dynamic programming on the block-cut
tree}\label{dynamic-programming-on-the-block-cut-tree}

Choose a root block \(B_0\) and orient the block-cut tree away from it.

For a nonroot block \(B\), let \(c\) be its parent cutvertex. Its
subtree contains no other block incident with \(c\). Thus, in the
subtree, requiring \(c\) to be on the outer face requires an outer-face
incidence in \(B\) itself.

Process the block-cut tree bottom-up.

\begin{itemize}
\tightlist
\item
  At each vertex \(v\) of a block, use as \(w(v)\) the maximum of the
  already computed values of the child branches attached there.
\item
  For an ordinary biconnected block \(B\), compute its message toward
  its parent by the constrained weighted algorithm (25).
\item
  A bridge block has one face, so its message is simply \[
  \max\{1,\max_{v\in V(B)}w(v)\}.
  \qquad\text{(28)}
  \]
\item
  At the root block \(B_0\), use the unconstrained weighted algorithm
  (24), or (28) if it is a bridge.
\end{itemize}

By Lemma 3 and induction, this computes the optimum among embeddings
whose outer face is incident with \(B_0\).

Finally, try every block as \(B_0\). Every embedding of a connected
graph with an edge has some block incident with its outer face. Thus the
minimum over root blocks is the unrestricted optimum.

For a disconnected graph, solve each component separately and place the
resulting embeddings side by side. The answer is the maximum of the
component answers. Restriction to a component gives the matching lower
bound.

\begin{center}\rule{0.5\linewidth}{0.5pt}\end{center}

\subsubsection{9. Polynomial running time and
reconstruction}\label{polynomial-running-time-and-reconstruction}

Let \(N=|V(G)|+|E(G)|\).

For one rooted SPQR computation:

\begin{itemize}
\tightlist
\item
  there are \(O(N)\) nodes and total skeleton size \(O(N)\);
\item
  every profile has \(O(N)\) entries;
\item
  an \(R\)-node uses weighted shortest paths and table lookups;
\item
  an \(S\)-node uses entrywise maxima;
\item
  a \(P\)-node uses at most \(O(N^4)\) operations by (22).
\end{itemize}

Thus \(O(N^5)\) operations is a safe coarse bound for one rooted
computation.

Trying root edges contributes a factor at most \(N\). Computing messages
for all blocks for a fixed root block contributes another factor at most
\(N\), and trying all root blocks contributes one more. This gives the
advertised coarse bound \[
O(N^8).
\]

All relevant values have \(O(\log N)\) bits. In particular, a branch's
optimum is at most its number of edges. The profile and scheduling
values are also bounded by a constant multiple of \(N\).

For reconstruction, retain minimizing choices:

\begin{itemize}
\tightlist
\item
  the selected root block and root edge;
\item
  each \(R\)-node reflection;
\item
  child profile arguments and orientations;
\item
  each \(P\)-node's central child and arm assignment;
\item
  a nearest incident face at every cutvertex.
\end{itemize}

The SPQR choices produce the block embeddings. Lemma 3 then specifies
how to insert the branch embeddings. Hence the algorithm returns an
embedding, not merely the optimum value.

For multigraphs, parallel edges are already covered by the dipole
skeletons. A loop can be treated as an elementary one-edge block, with
the same value formula as (28). Isolated vertices have no effect.

\begin{center}\rule{0.5\linewidth}{0.5pt}\end{center}

\subsubsection{10. Conclusion and
dependencies}\label{conclusion-and-dependencies}

The algorithm is exact. Its two essential ingredients are:

\begin{enumerate}
\def\labelenumi{\arabic{enumi}.}
\tightlist
\item
  \textbf{Invariant interface lengths:} a pertinent component's shortest
  dual distance between its two boundary faces does not depend on its
  embedding.
\item
  \textbf{A polynomial ordering step:} at a parallel node, at most one
  child receives genuinely two-sided distances; all other children lie
  on two one-sided arms, each optimally sorted by
  \(F_i(\ell_i)-\ell_i\).
\end{enumerate}

The first makes child optimization independent at fixed skeletons. The
second removes the exponential permutation choice at parallel skeletons.
Vertex weights and Lemma 3 handle all cutvertices.

This proves the polynomial-time claim. The external structural
dependency is the standard SPQR decomposition theorem and its
completeness for planar embeddings. I have not independently established
whether this particular algorithm, or an equivalent one, is already
present in the literature; the argument above is an independent proof
rather than a literature-status assertion.

\end{document}
